\subsection{Logarithmic tangent slopes and fixed-reference wall limits}
\label{sec:logarithmic-wall-heights}

The following argument supplies the logarithmic slope estimate needed
in the reference point-height comparison. In particular, it is not
enough to bound the maximum tangent slope linearly in the Faltings
height. We use ordinary derivatives, rather than divided powers, and
integral theta translations rather than powers of a single section.

\paragraph{Integral theta translations.}
Let \(A/K\) be a principally polarized abelian variety of dimension
\(g\ge1\), with semistable reduction, and let \(\mathcal E\) be the Lie
lattice of its N\'eron identity component, with the Riemann metric
of the polarization. Write \(\deg_n\) for arithmetic degree divided
by the degree of the number field, and put
\[
 H=2+g+\max(0,-\deg_n\mathcal E).
\]
Replacing the Riemann norm by the fixed normalization used for
Faltings metrics changes degrees by \(O(g)\) and slopes by \(O(1)\).
We prove, with absolute constants,
\[
                 \mu_{\max}(\mathcal E)\ll\log H.       \tag{LW1}
\]
All field extensions below are harmless for normalized line degrees.
In particular the semistable Lie lattice commutes with base change
and the metrics at the extended embeddings are the original metrics.

Set \(m=6\lceil H\rceil,\ n=m^2,\ R=m^{2g}\). After a finite
extension choose a symmetric rigidified ample line \(L\) inducing
the principal polarization. Use the Moret--Bailly model with all
\(A[m^2]\) as prescribed sections. The precise references are
Pazuki, \emph{Theta height and Faltings height}, Bull. Soc. Math.
France 140 (2012), 19--49, Theorem 3.4(i), (iv), (v),
Propositions 2.5 and 5.1
(\texttt{https://www.numdam.org/item/10.24033/bsmf.2623.pdf}).
The smooth semiabelian model is open in the N\'eron model and
has the same identity tangent lattice. Its cubical extension
\(\mathcal L\) has normalized Haar \(L^2\) metrics on its pushforward.
The rank-one pushforward \(\pi_*\mathcal L\) has normalized degree
\(-h_F(A)/2-g\log(2\pi)/4\) in Pazuki's conventions.
Principalize this invertible module after a further finite field
extension, using its pulled-back lattice, and take an integral
generator \(s\). Thus
\[
 \mathop{\rm av}_{\sigma}\log\|s\|_{2,\sigma}
             \le \tfrac12\max(0,-\deg_n\mathcal E)+C g .
\]
Here and below \(\mathop{\rm av}_{\sigma}\) averages all embeddings
with their usual multiplicities. Proposition 5.1 is important
integrally: multiplication by \(m\) and the good translations by
\(A[m^2]\) give isometries
\([m]^*\mathcal L\simeq\mathcal L^n\) and
\(t_x^*\mathcal L^n\simeq\mathcal L^n\).
Taking the translates of the pullback of \(s\) for representatives
of \(A[m^2]/A[m]\) gives \(R\) integral sections \(s_1,\ldots,s_R\)
spanning the generic \(H^0(A,L^n)\). Every one has sup norm
\(M_\sigma=\|s\|_{\sup,\sigma}\), not \(M_\sigma^n\).
Only the integral sections near the identity and their spanning
over the field will be used.

\paragraph{A uniform bound for the mean log sup norm.}
Write \(\rho_\sigma\) for the shortest nonzero period in the
Riemann metric. In the principally polarized complex torus the
Euclidean covolume for this metric is one. A ball of radius
\(r_\sigma=\min(1,\rho_\sigma/3)\) centered at any point embeds
in the torus. In a holomorphic Gaussian frame centered there,
the point norm is the absolute value of the holomorphic
coefficient times \(\exp(-\pi\|z\|^2/2)\), with a unit frame
at the center. The ball mean-value inequality gives
\[
 M_\sigma\le e^{\pi/2}
       \frac{\sqrt{g!}}{\pi^{g/2}r_\sigma^g}\|s\|_{2,\sigma}.
\]
For the inverse period length use Gaudron--R\'emond,
\emph{Th\'eor\`eme des p\'eriodes et degr\'es minimaux
d'isog\'enies}, Comment. Math. Helv. 89 (2014), 343--403,
Proposition 3.6
(\texttt{https://lmbp.uca.fr/\string~gaudron/art11.pdf}).
It bounds \(\mathop{\rm av}\rho_\sigma^{-2}\) by
\(11\max(1,h(A),\log\deg_L A)\). Its modified height differs
from the stable Faltings height by \(O(g)\);
\(\deg_L A=g!\). Jensen's inequality for \(\log\) then gives
\(\mathop{\rm av}\log^+(\rho_\sigma^{-1})\ll\log H\).
Consequently
\[
 \mathop{\rm av}_{\sigma}\log M_\sigma
                \le H/2+C g\log H.                    \tag{LW2}
\]
All constants are independent of \(m\), the field extension and the
abelian variety.

\paragraph{Integral directional jets.}
Fix a nonzero integral tangent vector \(v\).
Equip the coefficient module \(\mathcal O_K^R\) with the standard
Hermitian norm. For \(j\ge0\), let \(F_j\) be the saturated kernel
of the first \(j\) derivatives at the identity of
\(\sum_i a_i s_i\) along the invariant vector field \(v\).
On \(F_j\) the derivative of order \(j\) defines an integral
map to the identity fibre, identified with \(\mathcal O_K\)
using the rigidification of \(\mathcal L^n\). To check this
assertion use, at each finite place, a local algebraic frame
with the prescribed value at one. The invariant derivation
preserves the integral local ring. Changing such a frame
does not change the derivative on \(F_j\), since every lower
derivative of the section is zero. The same argument identifies
the archimedean derivative with the derivative in the holomorphic
Gaussian frame on the line \(zv\) in the universal cover. No
factor \(1/j!\) has been asserted integral.

Put \(J=n/3\), an integer. For every \(0\le j\le J\), the
derivative is nonzero on \(F_j\otimes K\). Indeed \(L^3\) is
very ample; take a section \(a\) nonzero at one and a section
\(b\) with simple zero along the specified tangent direction.
Then \(a^{J-j}b^j\) has exactly order \(j\) in that
direction. The translated sections span its coefficient space
over \(K\), and clearing denominators gives a coefficient
vector in \(F_j\). Thus \(F_j/F_{j+1}\) has rank one.
The argument also works for \(g=1\).

For a coefficient vector of norm at most one, the holomorphic
coefficient on \(zv\) is bounded by
\(\sqrt R M_\sigma\exp(\pi n\|v_\sigma\|^2 |z|^2/2)\).
Cauchy's formula at radius \(\sqrt{j/(\pi n)}/\|v_\sigma\|\)
and \(j!\le j^j\) show that the norm of the integral jet map
on \(F_j/F_{j+1}\), with the quotient metric, is at most
\[
 B_{\sigma,j}=\sqrt R\, M_\sigma
         (\sqrt{\pi e n j}\,\|v_\sigma\|)^j
 \quad (j\ge1),\qquad B_{\sigma,0}=\sqrt R\, M_\sigma .
\]
Taking the infimum over lifts is legitimate because the jet
annihilates \(F_{j+1}\). An integral generically injective map
of a Hermitian line into the trivial line of norm one bounds
its normalized degree above by the mean log operator norm.
On the other hand the degree of the quotient
\(\mathcal O_K^R/F_{J+1}\) is nonnegative: its dual is a
saturated submodule of a standard Hermitian lattice, whose
determinant has degree at most zero by the product formula
on any nonzero Pl\"ucker coordinate. Additivity with the
subquotient metrics now gives
\[
 0\le\sum_{j=0}^J\deg_n(F_j/F_{j+1})
 \le (J+1)(\tfrac12\log R+\mathop{\rm av}\log M_\sigma)
 +\frac{J(J+1)}2
       (\tfrac12\log(\pi e n J)+\mathop{\rm av}\log\|v_\sigma\|).
\]
In particular
\[
 -\mathop{\rm av}\log\|v_\sigma\|
 \le \tfrac12\log(\pi e n J)
       +\frac{\log R+2\mathop{\rm av}\log M_\sigma}{J}
 \ll\log H.                                           \tag{LW3}
\]
For the last estimate \(J=m^2/3\ge12H^2,\ 
\log R=2g\log m,\ g\le H\), and (LW2) suffice.
This bound applies to every integral vector, also after further
field extension. For any saturated line after finite extension,
principalize it and take a generator; its normalized degree is
bounded by (LW3). The absolute Siegel lemma (Zhang; see
Gaudron--R\'emond, \emph{Minima, pentes et alg\`ebre
tensorielle}, Israel J. Math. 195 (2013), 565--591,
Theorem 1.1) bounds the maximum slope by the supremum of such
absolute line degrees plus \(\log g\). This proves (LW1).
Rank zero requires no assertion about a maximum slope.

\paragraph{The semiabelian and flag comparisons.}
For the boundary application, the following elementary
extension observation avoids any surjectivity assertion about
the integral map onto the abelian quotient. Suppose that the
actual tangent module \(E\) has a saturated toric kernel \(T\),
and its torsion-free quotient injects as a full lattice into
the abelian tangent module \(E_A\). Assume the toric metric
is comparable to a standard cocharacter lattice with fixed
constants, and the map \(E/T\to E_A\) with its quotient metric
has uniformly bounded archimedean operator norms. These
are metric requirements, not just generic exactness.
Writing \(C\) for a common constant (for a fixed tower),
\[
 -\deg_n E_A\le-\deg_n E+C,\qquad
 \mu_{\max}(E)\le
       \max(\mu_{\max}(T),\mu_{\max}(E_A)+C).
\]
For the first inequality use the determinant of the integral
injection, whose finite index has the favorable sign.
For the second intersect any saturated subbundle with \(T\)
and take its image in \(E_A\); its quotient metric dominates
the induced metric up to the stated operator-norm constant.
Add degrees and divide by the sum of the ranks.
This also proves the assertion when one block is zero.
In particular (LW1) gives a logarithmic slope bound on
\(E\) for the reference-point argument under these exact
uniform metric identifications.

Here is the required flag-height consequence, stated with
the geometric identifications explicit. On a fixed projective
reference flag variety let \(h_A\ge0\) be a height of the fixed
ample auxiliary line (a bounded shift is permitted).
Assume the actual orthogonal flag comparison of
\texttt{FibreInput}:
the first-block root heights differ by uniformly bounded
amounts from the graded degrees of the saturated flag \(W\);
the last block differs by bounded amounts from the negatives
of the graded degrees of the reverse flag \(S\).
The total signed determinant is bounded, with its original
flat unit; the paired orthogonal determinants give
\(-\deg_n E\ll1+h_A\). The latter bound also follows from
the usual comparison of the heights of the fixed determinant
lines with the ample height. Retain the preceding toric and
abelian metric comparison. For dominant weights \(u\) with
\(\|u\|_\infty\le2\), summation by parts writes the negative
weighted root sum as a sum of nonnegative adjacent gaps
times negative degrees of \(W\)-prefixes and \(S\)-prefixes,
including the middle gap times \(-\deg_n W\), plus a bounded
flat-determinant term. Every prefix degree is at most
\(N\mu_{\max}(E)\) when \(\mu_{\max}(E)\ge0\);
one may use its positive part otherwise.
Hence (LW1) gives
\[
       h_{\mathrm{ref}}(u,x)\ge-C\log(2+h_A(x)).        \tag{LW4}
\]
The hypotheses here refer to the actual reference point
metrics, including at its boundary. The new slope proof
does not identify a generically isomorphic lattice with
the prescribed Hermitian fibre.

\paragraph{Wall convergence without a uniform smoothing rate.}
We record a strengthened comparison on the linked tower
of Section~\ref{sec:flag-finite-wall}. The following proof uses
the actual model nefness and the source smoothing
normalization, and needs convergence of mixed intersections
only on one fixed reference model.

More precisely, fix a tower. Retain the established common
model and isometric linkage, with its normalized cover
degrees, and the actual auxiliary ample line with its fixed
generic metric. Assume the reference flag and quotient
metric identifications in (LW4), and the original prefix
smoothing estimates: for dominant weights of bounded norm
the smoothed height minus the fixed smooth reference height
is between \(-C(1+\log(1/s))\) and \(C\) at parameter
\(t=\exp(-1/s^2)\), before the finite-place changes.
These are the bounds obtained by summation by parts in
\texttt{RootTowerUniformHeightBound} from the prefix potentials
and the bounded final flat prefix. Also retain the original
semipositive smoothings, the fibre nefness threshold
\texttt{ModelNef}, \texttt{M2}, the nef auxiliary line,
the top-degree proportionality, and the qualitative limits
of all degree-\(d+1\) mixed intersections on the one reference
model. In the common models the finite-place sandwiches
are the original root-lattice sandwiches, not bounds on
heights of individually chosen points.

We first derive the fixed-tower coefficient bound. Choose
finitely many fixed rational strictly dominant weights
and strictly positive rational auxiliary coefficients
so their \((d+1)\)-st powers span the polynomial space of
degree \(d+1\) in the root and auxiliary variables. This
is possible by the polynomial identity principle on an
open cone. Choose a fixed cutoff \(B_0\) so that all these
lines are nef on the fibres at primes above \(B_0\), by
the original \texttt{M2} threshold. Use a common comparison
model for \(B_0,B\). For each of these fixed lines, the
difference of the pullbacks is vertical and bounded in
both directions, at each new prime \(\ell\), by a constant
times \(\gamma_\ell\operatorname{div}(\ell)\).
Here the original corrected-lattice sandwich gives
\(\gamma_\ell\ll1/(\ell-1)\) for odd new primes; the finitely
many exceptional primes contribute a fixed constant.
This sandwich passes to the saturated flag subquotients:
intersect the two bounding lattices with a flag subspace
and take images in the next quotient. Common model
pullbacks use the already locally free flag bundles.
At every new fibre both line bundles are nef.
Telescoping their top arithmetic powers and intersecting
the two divisor bounds against the mixed nef products
bounds the absolute change by a constant times
\(\sum_{\ell\le B}\gamma_\ell\log\ell\).
The mixed products against \(\operatorname{div}(\ell)\)
have the corresponding generic intersection degrees after
normalization by the cover degree; those degrees are fixed.
There is no archimedean difference at a common smoothing
parameter by isometric linkage. Chebyshev's prime bound
and partial summation give \(O(1+\log B)\) for the sum.
Passing to the reference limits and interpolating in the
fixed polynomial space therefore gives
\[
               |\operatorname{PP}_B(F)|\ll_F1+\log B.   \tag{LW5}
\]
Equivalently this holds with a fixed constant for a fixed
monomial basis. No uniform bound in the varying tower
dimension is needed by the finite wall test argument.

In fact the difference of the normalized arithmetic
intersection functionals at two smoothing parameters is
independent of \(B\). Pull back both models to the common
model; their metrics agree at each parameter, and the
change of an arithmetic intersection with the metric is
an archimedean integral. The cover degree cancels.
Hence
\[
 \operatorname{Ar}_{B,t}(F)-\operatorname{PP}_B(F)
   =\operatorname{Ar}_{B_0,t}(F)-\operatorname{PP}_{B_0}(F).
                                                               \tag{LW6}
\]
Existence of the limits on the right gives existence on
the left, and the right side is \(o(1)\) uniformly over
bounded coefficient vectors \(F\), by finite dimensionality.
This removes the need for a quantitative smoothing rate
uniform in the cutoff. It does not establish the reference
mixed-intersection limits themselves.

Take one of the rational wall tests \(u\) of norm less than
\(1/2\), with exactly the selected cross-block adjacent gap
zero and all other adjacent gaps strict. Put \(\theta=5f/B\),
raise the first block in that label by \(2\theta\), obtaining
\(u_\theta\), and put \(s=\theta^2,\ v=u_\theta-s\nu\)
with the original fixed \(\nu\).
For large \(B\) the other gaps remain strictly positive and
\((B-2)\) times each cross gap exceeds \(4f\|v\|_\infty\).
Thus the original fibre thresholds and the smoothing
positivity make the rational smoothed line
\(\overline D_B=\overline L_B(v)+s\overline A_B\)
relatively nef, with ample generic fibre (the strictly
dominant class is nef and the auxiliary line is ample).
Clear denominators when applying arithmetic intersection
and successive-minima inequalities. Its geometric degree
is \(O(\theta)\): the root-only degree vanishes at \(u\)
by the Vandermonde product and changes by \(O(\theta)\),
and adding the auxiliary class changes the degree by
\(O(s)\). These polynomial estimates use the fixed
normalized generic intersection functional.

By (LW4), the original prefix bounds and the same
finite-place sandwich, all points satisfy
\[
 h_{\overline D_B}(x)\ge s h_A(x)-C\log(2+h_A(x))
                       -C\log(B/s)-C.
\]
Minimizing the first two terms for \(h_A\ge0\) gives an
all-point lower bound \(-C'\log(B/s)\). The elementary
inequality follows either by differentiation or from
\(\log z\le az-1-\log a\) for \(z,a>0\).
For the upper bound there is a Zariski dense family of
points with bounded \(h_A\) on the fixed reference variety:
a finite projective linear projection to projective
space gives such points above the dense set of torsion
points in its affine torus. All lifts have bounded
height by the pullback ample height identity. Use the
dominant common covers to lift this dense set to each
tower model, discarding any closed exceptional locus.
Their normalized reference heights stay bounded.
The upper prefix comparison and the finite sandwiches
give an upper bound \(C'\log(B/s)\) on that dense set.

The arithmetic successive-minima theorem (Zhang,
\emph{Positive line bundles on arithmetic varieties},
J. Amer. Math. Soc. 8 (1995), 187--221, (5.2)) sandwiches
the top arithmetic power divided by \((d+1)\) times the
generic degree between the all-point infimum and the
essential minimum for this relatively semipositive line.
The preceding lower and dense upper bounds therefore give
\[
             |\operatorname{Ar}_{B,t}(D_B^{d+1})|
                    \ll\theta\log(B/s)=O(\log B/B).
\]
Use (LW6) on the bounded coefficient family of these
powers and then (LW5) on the difference from \(\ell_u^{d+1}\),
whose coefficients are \(O(\theta)\). This proves
\[
                  \operatorname{PP}_B(\ell_u^{d+1})\longrightarrow0.
                                                               \tag{LW7}
\]
The argument applies separately to every fixed wall test,
as required for the exact cutoff identity. Unlike the
original uniform \texttt{ThmW}, this statement needs no
estimate for the rate of the reference smoothing limit.
The qualitative reference convergence, semipositive
metrics, actual uniform fibre identifications and
proportionality are retained hypotheses of this
application; the argument has not supplied them simply
from equality of generic line bundles.

\paragraph{Base-ideal multiplicities along transverse arcs.}
There is a separate way to use the already reconstructed
thin-region lower Bergman estimate for the finite divisor
certificate. On the normalized blow-up of the actual
cusp-section system, let \(D\) be a prime divisor of its
base divisor \(E\), and take a transverse complex analytic
arc through a general smooth point of \(D\), with punctured
arc in the interior. Let \(b_D=\operatorname{ord}_D E\).
For a fixed orthonormal basis of the complex space of
cusp sections, write its coefficients in a local frame of
the pulled-back weight line. Their common minimum order
is exactly \(b_D\); after division by \(t^{b_D}\) at least
one coefficient has nonzero limit. Thus if the metric
has sublogarithmic growth in this frame along the arc,
\[
       \log\mathcal B(x(t))=-2 b_D\log(1/|t|)
                                  +o(\log(1/|t|)).
\]
This uses the full actual section system, not an arbitrary
nonzero section, and applies to exceptional divisors as
well as strict transforms. At a point of the interior
where the system is basepoint free there is no such
divisor. The source lower bound
\(\mathcal B(x)\ge\exp(-\Phi(x)/625)\) then gives
\[
 b_D\le\frac1{1250}
             \limsup_{t\to0}\frac{\Phi(x(t))}{\log(1/|t|)}.
                                                               \tag{LW8}
\]
In the source divisor construction the desired theta
section has weight \(8N\,\mathrm{hodgeWt}\), and lower
valuation \(8\phi_D/625\) at the original monodromy
valuation of \(D\). If this section is constructed with
that lower bound and if the arc comparison gives
\(\limsup\Phi(x(t))/\log(1/|t|)\le6\phi_D\), then (LW8)
implies \(8b_D\le24\phi_D/625\).
Its cube has the weight \(8(A-A')\) and dominates
\(8E\), proving (FW1). This is a sufficient finite
arc estimate, not a proof that the theta section or
the stated monodromy comparison has already been supplied.
Both must use the actual cusp system, canonical metric
and normalized divisorial valuation.
